\section{Related work and what is new here}
\label{sec:related}

\paragraph{Calibration error in adaptive testing.} Treating estimated item parameters
as true inflates measurement precision and biases ability estimates, and adaptive
selection amplifies the problem because it seeks the items that look most informative,
the \emph{capitalization on calibration error}
\citep{vanderlinden2000capitalization,olea2012capitalization,patton2013influence}.
Remedies propagate parameter uncertainty into scale scores
\citep{zhang2011investigating,yang2012characterizing}, replace the item response
function by its expectation over the parameter posterior---the expected response
function \citep{mislevy1994erf}, which we use---or treat all parameters in a fully
Bayesian way. \citet{vanderlinden2020fast} give a fast Markov chain Monte Carlo
algorithm for Bayesian adaptive testing that samples the joint posterior of ability and
item parameters during the test, using stored posterior draws of the item parameters
from a Bayesian calibration. \citet{fink2025accounting} compare corrections for
calibration error in simulated 2PL tests with 100--500 calibration responses per item
and find the fully Bayesian approach, with item selection by Bayesian maximum
information averaged over ability and item draws, the most accurate. We therefore do
\emph{not} claim to be the first to integrate calibration uncertainty into adaptive
testing. Lemma~\ref{lem:exact} shows that our calibration-marginal model is exactly
the fully Bayesian model with the calibrators' abilities and the hyperprior held at point
estimates; study~2 compares the two directly.

\paragraph{Item selection.} Classical CAT selects the item with maximum Fisher
information at the current ability estimate \citep{lord1980applications}; Bayesian
variants weight information by the ability posterior
\citep{owen1975bayesian,vanderlinden1998bayesian}, use Kullback--Leibler information
\citep{chang1996global} or other criteria \citep{veerkamp1997some}, and differ mainly
early in a test \citep{chen2000comparison}; randomesque selection among the most
informative items limits exposure \citep{kingsbury1989procedures}. Mutual-information
selection was proposed for classification testing \citep{weissman2007mutual} and
multidimensional CAT \citep{wang2011item}; in experimental design it is Lindley's
information \citep{lindley1956measure,chaloner1995bayesian,rainforth2024modern} and in
active learning BALD \citep{houlsby2011bald}, whose motivation---entropy conflates
epistemic and aleatoric uncertainty---is Proposition~\ref{prop:decomp}. Uncertainty
sampling \citep{lewis1994sequential} and expected error reduction
\citep{roy2001toward} are the active-learning counterparts of our entropy and
score-targeted criteria.

\paragraph{IRT and reduced evaluation for LLMs.} IRT has been used to build evaluation
scales \citep{lalor2016building}, to weight leaderboard examples
\citep{rodriguez2021evaluation}, to select anchor points and compressed benchmarks
\citep{vivek2024anchor,polo2024tinybenchmarks,kipnis2025metabench} and to amortize
calibration with content-based difficulty predictors \citep{truong2025reliable}.
tinyBenchmarks \citep{polo2024tinybenchmarks} choose about 100 anchor items by clustering
item embeddings from hundreds of training models and combine the anchors with IRT
predictions (gp-IRT). ATLAS \citep{li2026atlas} runs Fisher-information CAT with a 3PL
calibrated on thousands of leaderboard models and reports up to 90\% fewer items;
Fluid Benchmarking \citep{hofmann2025fluid} does the same with about one hundred
calibration models. MINCE \citep{das2026mince} uses a few calibration models only to size
a random subset by Monte Carlo, without an IRT model. Other work extends adaptive testing
to continuous scores \citep{balkir2026confident}, accounts for prompt sensitivity with
parameters treated as known \citep{jo2025measure}, links new datasets by fixed-parameter
calibration \citep{habba2026growing}, uses neural or medical item banks
\citep{zhou2026lost,zheng2026medcat}, or applies a Laplace approximation to neural IRT
difficulties with selection by posterior-expected Fisher information
\citep{mandujano2026laplace}; a position paper argues for psychometric testing of AI
systems \citep{zhuang2025position}, and a simulation study finds item-level IRT
inferences unreliable when few calibration models are available
\citep{jiang2026trust}---the regime of this paper. Error bars for LLM evaluation are
studied by \citet{madaan2024quantifying} and \citet{miller2024adding}.

\paragraph{What is new.} The ingredients---expected response functions, Bayesian
adaptive testing, information-based selection---are classical. What this paper adds is
(i) their systematic study in the regime that LLM evaluation actually faces, with
calibration panels of 4 to 64 systems instead of hundreds to thousands of examinees, on
eight real benchmarks from five domains and with independent targets;
(ii) theory for that regime: exactness of the calibration-marginal model given the
calibrators (Lemma~\ref{lem:exact}), an information cap that holds exactly for any link
under a Gaussian difficulty posterior (Proposition~\ref{prop:cap}), the relation between entropy, Fisher information and EIG
under calibration uncertainty, and a one-step score-targeted criterion; (iii) head-to-head
comparisons, under analysis plans frozen before execution, with the competitors that
matter in this regime---a regularized Rasch-type test, Bayesian adaptive tests that also
integrate the calibrators' abilities,
ATLAS-style 3PL CAT, tinyBenchmarks, MINCE, and uniform subsampling with and without the
same calibration; and (iv) an honest account of uncertainty quantification, including a
leave-one-calibrator-out recalibration of score intervals and a cost analysis.

\paragraph{Self-evaluation and circularity.} Using models to judge models raises
circularity concerns \citep{zheng2023judging}. The framework allows models to generate
items but requires an independent oracle for verdicts; this aspect was explored only in
a historical pilot (Section~\ref{sec:pilot}) and is not part of the evidence of this
paper.
